\chapter{Data Understanding I: Data Types, Attributes, and Descriptive Statistics}
\label{ch:data-understanding-types-statistics-en}

\section{Understanding Data and Dataset Taxonomy}

\textbf{Data Understanding} is the diagnostic investigation in which the analyst examines data quality, understands domain feature semantics, and uncovers initial indications of noise, extreme outliers, or missing values.

\begin{noteblock}{Core Objectives of Data Understanding}
\begin{itemize}
    \item What types of attributes constitute the dataset?
    \item What is the overall data quality, and what syntactic/semantic errors exist?
    \item Are there univariate or multivariate outliers present?
    \item What pairwise correlations and dependencies connect features?
    \item How are missing values distributed, and what mechanism generated them?
    \item Is feature transformation or feature engineering required prior to mining?
\end{itemize}
\end{noteblock}

\subsection{Fundamental Structure: Objects and Attributes}
A dataset is a formal collection of **data objects** (also termed *records, points, samples, entities, instances, or cases*) described by a set of **attributes** (known as *variables, features, fields, or dimensions*).
\begin{itemize}
    \item An \textbf{object} represents a single real-world entity (a customer, a bank transaction, a sensor reading).
    \item An \textbf{attribute} denotes a measurable property or characteristic of that entity (height, price, temperature, category).
\end{itemize}

\subsection{Classification of Data Structures}
Datasets differ fundamentally by their underlying structural topology:

\begin{figure}[H]
\centering
\begin{tikzpicture}[
    catNode/.style={rectangle, rounded corners=6pt, draw=navy!80!black, fill=navy!6, thick, text width=3.5cm, minimum height=3.0cm, align=left, font=\footnotesize, inner sep=6pt},
    arrowTax/.style={-Stealth, thick, draw=navy!85!black}
]
    \node[rectangle, rounded corners=6pt, draw=purple!85!black, fill=purple!12, very thick, font=\small\bfseries, align=center, text width=6.6cm, inner sep=7pt] (root) at (0, 2.5) {DATASET TAXONOMY};

    \node[catNode, anchor=north] (rec) at (-4.7, 1.1) {
        \textbf{Record Data}\\[3pt]
        $\bullet$ Matrix $m \times n$\\
        $\bullet$ Text (BoW)\\
        $\bullet$ Transactions
    };

    \node[catNode, anchor=north] (graph) at (0, 1.1) {
        \textbf{Graph Data}\\[3pt]
        $\bullet$ Web Hyperlinks\\
        $\bullet$ Molecular Graphs\\
        $\bullet$ Social Networks
    };

    \node[catNode, anchor=north] (ord) at (4.7, 1.1) {
        \textbf{Ordered Data}\\[3pt]
        $\bullet$ Time Series\\
        $\bullet$ Sequences\\
        $\bullet$ Genetic Sequences\\
        $\bullet$ Spatio-Temporal
    };

    \coordinate (split) at (0, 1.6);
    \draw[thick, draw=navy!85!black] (root.south) -- (split);
    \draw[arrowTax] (split) -| (rec.north);
    \draw[arrowTax] (split) -- (graph.north);
    \draw[arrowTax] (split) -| (ord.north);
\end{tikzpicture}
\caption{Structural taxonomy of major dataset formats.}
\label{fig:dataset-taxonomy-en}
\end{figure}

\subsubsection{1. Record Data}
\begin{itemize}
    \item \textbf{Data Matrix:} When all $m$ objects share the same fixed set of $n$ continuous numerical attributes, the dataset is modeled as an $m \times n$ matrix:
    \begin{equation}
        D = \begin{bmatrix}
        x_{11} & x_{12} & \cdots & x_{1n} \\
        x_{21} & x_{22} & \cdots & x_{2n} \\
        \vdots & \vdots & \ddots & \vdots \\
        x_{m1} & x_{m2} & \cdots & x_{mn}
        \end{bmatrix} \in \mathbb{R}^{m \times n}
    \end{equation}
    where each row $i$ is a point in continuous Euclidean space $\mathbb{R}^n$, and each column is a coordinate axis.
    \item \textbf{Document Data (Document-Term Matrix):} Unstructured texts modeled via \textit{Bag-of-Words}. Columns represent vocabulary terms and values record word frequencies:
    \begin{table}[H]
    \centering
    \resizebox{\linewidth}{!}{%
    \begin{tabular}{l *{10}{c}}
    \toprule
    \textbf{Document} & \textbf{team} & \textbf{coach} & \textbf{play} & \textbf{ball} & \textbf{score} & \textbf{game} & \textbf{win} & \textbf{lost} & \textbf{timeout} & \textbf{season} \\
    \midrule
    Doc 1 & 3 & 0 & 5 & 0 & 2 & 6 & 0 & 2 & 0 & 2 \\
    Doc 2 & 0 & 7 & 0 & 2 & 1 & 0 & 0 & 3 & 0 & 0 \\
    Doc 3 & 0 & 1 & 0 & 0 & 1 & 2 & 2 & 0 & 3 & 0 \\
    \bottomrule
    \end{tabular}%
    }
    \caption{Document-Term Matrix representation for text collections.}
    \label{tab:doc-term-matrix-en}
    \end{table}
    \item \textbf{Transaction Data:} Each record is identified by a unique TID and contains a variable-length set of discrete items (e.g., supermarket market basket).
\end{itemize}

\subsubsection{2. Graph Data}
Explicitly models pairwise entities: World Wide Web hyperlinks, molecular structures (e.g., benzene $C_6H_6$ covalent bonds), and citation networks.

\subsubsection{3. Ordered Data}
Data where temporal or spatial precedence governs informational semantics:
\begin{itemize}
    \item \textbf{Time Series:} Continuous or discrete time-indexed readings with high autocorrelation.
    \item \textbf{Sequential Data:} Chronologically ordered event sequences (e.g., $(A, B) \to (D) \to (C, E)$).
    \item \textbf{Genomic Sequences:} Nucleotide polymer chains (`A`, `C`, `G`, `T`).
    \item \textbf{Spatio-Temporal Data:} Space-time indexed matrices (e.g., global sea-surface temperature grids).
\end{itemize}

\FloatBarrier
\section{Attribute Types and Measurement Scales}

Mathematical admissibility of operations on attribute domains:
\begin{itemize}
    \item \textbf{Distinctness ($=$ or $\neq$):} Testing equality between values.
    \item \textbf{Order ($<$ or $>$):} Natural total or partial ordering.
    \item \textbf{Difference ($+$ or $-$):} Constant-metric measurable intervals.
    \item \textbf{Ratio ($*$ or $/$):} Natural absolute zero (total absence of measured property).
\end{itemize}

\begin{table}[H]
\centering
\small
\begin{tabularx}{\linewidth}{l c c c c Y}
\toprule
\textbf{Scale} & \textbf{Distinct.} & \textbf{Order} & \textbf{Diff.} & \textbf{Ratio} & \textbf{Permissible Statistical Operators} \\
\midrule
\textbf{Nominal}  & \checkmark & --         & --         & --         & Mode, entropy, $\chi^2$ test \\
\textbf{Ordinal}  & \checkmark & \checkmark & --         & --         & Median, quantiles, Spearman rank $\rho$ \\
\textbf{Interval} & \checkmark & \checkmark & \checkmark & --         & Mean, standard deviation, Pearson $r$ \\
\textbf{Ratio}    & \checkmark & \checkmark & \checkmark & \checkmark & Geometric mean, coefficient of variation \\
\bottomrule
\end{tabularx}
\caption{Hierarchy and mathematical properties of measurement scales.}
\label{tab:measurement-scales-en}
\end{table}

\subsection{Categorical (Qualitative) Attributes}
\begin{enumerate}
    \item \textbf{Nominal:} Allows only equality comparison ($=, \neq$). Examples: eye color, zip code, tax ID.
    \item \textbf{Binary:} Nominal attributes with exactly two values $\{0, 1\}$.
    \begin{itemize}
        \item \textit{Symmetric Binary:} Both outcomes carry equal weight (e.g., gender Male/Female).
        \item \textit{Asymmetric Binary:} One outcome is rare or of focal interest ($1 = \text{positive medical test}$, $0 = \text{negative}$).
    \end{itemize}
    \item \textbf{Ordinal:} Natural ordering ($<, >$), but step distance is non-constant. Examples: Mohs hardness scale, ratings $\{\text{poor}, \text{fair}, \text{good}\}$, letter grades $\{A, B, C, D, F\}$.
\end{enumerate}

\subsection{Numeric (Quantitative) Attributes}
\begin{enumerate}
    \item \textbf{Interval-Scaled:} Differences are meaningful, but the scale **lacks an absolute zero**. Examples: Celsius or Fahrenheit temperature ($20^\circ\text{C}$ is not ``twice as hot'' as $10^\circ\text{C}$, since in Kelvin $293.15\,\text{K} \neq 2 \cdot 283.15\,\text{K}$); calendar dates.
    \item \textbf{Ratio-Scaled:} Possesses a **natural absolute zero**. Examples: length, weight, age, Kelvin temperature ($0\,\text{K}$ absolute zero), monetary counts. Stating that a 3-hour match is 50\% longer than a 2-hour match ($3/2 = 1.5$) is mathematically valid.
\end{enumerate}

\subsection{Discrete vs Continuous Variables}
\begin{itemize}
    \item \textbf{Discrete Attribute:} Finite or countably infinite set of values, represented by integers.
    \item \textbf{Continuous Attribute:} Real-valued domain ($\mathbb{R}$), represented by floating-point numbers.
\end{itemize}

\FloatBarrier
\section{Data Quality Issues}

Flawed data invalidates downstream analytical pipelines (*Garbage In, Garbage Out*). Redman (2004) estimated poor data quality costs organizations $10\%\div 20\%$ of total revenue.

\begin{figure}[H]
\centering
\begin{tikzpicture}[
    dimCard/.style={rectangle, rounded corners=6pt, draw=navy!80!black, fill=navy!6, thick, text width=6.0cm, minimum height=2.8cm, align=left, font=\footnotesize, inner sep=6pt},
    arrowStyle/.style={-Stealth, thick, draw=navy!85!black}
]
    \node[rectangle, rounded corners=6pt, draw=purple!85!black, fill=purple!12, very thick, font=\small\bfseries, align=center, text width=8.6cm, inner sep=6pt] (root) at (0, 2.7) {
        DATA QUALITY DIMENSIONS\\[2pt]
        \footnotesize \normalfont Critical factors for analytical reliability
    };

    % Left Column
    \node[dimCard, anchor=north] (leftCol) at (-3.6, 1.1) {
        \textbf{Accuracy and Completeness}\\[4pt]
        $\bullet$ \textbf{Syntactic}: Domain violations\\[3pt]
        $\bullet$ \textbf{Semantic}: Factually untrue values\\[3pt]
        $\bullet$ \textbf{Completeness}: Null values (\texttt{NaN})
    };

    % Right Column
    \node[dimCard, anchor=north] (rightCol) at (3.6, 1.1) {
        \textbf{Dynamics and Integrity}\\[4pt]
        $\bullet$ \textbf{Class Imbalance}: Rare classes\\[3pt]
        $\bullet$ \textbf{Timeliness}: Stale or obsolete data\\[3pt]
        $\bullet$ \textbf{Uniqueness}: Duplicate records
    };

    \coordinate (split) at (0, 1.65);
    \draw[thick, draw=navy!85!black] (root.south) -- (split);
    \draw[arrowStyle] (split) -| (leftCol.north);
    \draw[arrowStyle] (split) -| (rightCol.north);
\end{tikzpicture}
\caption{Core dimensions of data quality.}
\label{fig:data-quality-en}
\end{figure}

\FloatBarrier
\section{Univariate Distribution Inspection}

\subsection{Skewness and Multimodality}
The shape of continuous probability density functions reveals underlying mechanisms:

\begin{figure}[H]
\centering
\begin{tikzpicture}[scale=0.85]
    % 1. Symmetric
    \begin{scope}[xshift=-5.5cm]
        \node[font=\footnotesize\bfseries, text=blue!90!black] at (0, 2.8) {Symmetric (Gaussian)};
        \node[font=\tiny, text=gray!80!black] at (0, 2.4) {Mean = Median = Mode};
        \draw[-Stealth, draw=gray!70] (-2.2, 0) -- (2.2, 0);
        \draw[very thick, draw=blue!80!black] (-2, 0) .. controls (-1, 0.1) and (-0.8, 1.8) .. (0, 2.0)
                                                    .. controls (0.8, 1.8) and (1, 0.1) .. (2, 0);
        \draw[dashed, red!80!black] (0, 0) -- (0, 2.0);
    \end{scope}

    % 2. Right-Skewed
    \begin{scope}[xshift=0cm]
        \node[font=\footnotesize\bfseries, text=orange!90!black] at (0, 2.8) {Right-Skewed (Positive)};
        \node[font=\tiny, text=gray!80!black] at (0, 2.4) {Mode $<$ Median $<$ Mean};
        \draw[-Stealth, draw=gray!70] (-2.2, 0) -- (2.2, 0);
        \draw[very thick, draw=orange!85!black] (-2, 0) .. controls (-1.6, 0.3) and (-1.2, 1.9) .. (-0.8, 2.0)
                                                        .. controls (-0.4, 2.1) and (0.2, 0.7) .. (1.0, 0.3)
                                                        .. controls (1.5, 0.1) and (1.9, 0.05) .. (2.2, 0);
        \draw[dashed, red!80!black] (-0.8, 0) -- (-0.8, 2.0) node[below, font=\tiny] {Mode};
        \draw[dashed, teal!80!black] (-0.3, 0) -- (-0.3, 1.4) node[below, font=\tiny] {Med.};
        \draw[dashed, purple!80!black] (0.4, 0) -- (0.4, 0.6) node[below, font=\tiny] {Mean};
    \end{scope}

    % 3. Bimodal
    \begin{scope}[xshift=5.5cm]
        \node[font=\footnotesize\bfseries, text=teal!90!black] at (0, 2.8) {Bimodal (Two Peaks)};
        \node[font=\tiny, text=gray!80!black] at (0, 2.4) {Latent subpopulations};
        \draw[-Stealth, draw=gray!70] (-2.2, 0) -- (2.2, 0);
        \draw[very thick, draw=teal!85!black] (-2, 0) .. controls (-1.5, 0.2) and (-1.3, 1.8) .. (-1.0, 1.8)
                                                      .. controls (-0.7, 1.8) and (-0.4, 0.6) .. (0, 0.6)
                                                      .. controls (0.4, 0.6) and (0.7, 1.8) .. (1.0, 1.8)
                                                      .. controls (1.3, 1.8) and (1.5, 0.2) .. (2, 0);
    \end{scope}
\end{tikzpicture}
\caption{Univariate distribution profiles: symmetry, positive skewness, and multimodality.}
\label{fig:distribution-shapes-en}
\end{figure}

\begin{itemize}
    \item \textbf{Symmetric Distribution:} Reflectional symmetry about the center: $\text{Mean} = \text{Median} = \text{Mode}$.
    \item \textbf{Positively Skewed (Right-Skewed):} Long tail towards higher values (e.g., individual wealth):
    \begin{equation}
        \text{Mode} < \text{Median} < \text{Mean}
    \end{equation}
    \item \textbf{Negatively Skewed (Left-Skewed):} Tail stretching towards small values:
    \begin{equation}
        \text{Mean} < \text{Median} < \text{Mode}
    \end{equation}
    \item \textbf{Bimodal / Multimodal Distribution:} Signals overlapping subpopulations (e.g., adult height without conditioning on biological sex).
\end{itemize}

\subsection{Visual Tools: Bar Charts vs Histograms}
\begin{itemize}
    \item \textbf{Bar Chart:} Plots discrete categories with visible spacing between bars.
    \item \textbf{Histogram:} Partitions continuous domains into contiguous intervals (*bins*); rectangle area is proportional to empirical sample counts.
\end{itemize}

\paragraph{Bin Sizing and Sturges' Rule:}
Under approximately normal distributions, Sturges' rule determines optimal bin counts $k$:
\begin{equation}
    k = \lceil \log_2(n) + 1 \rceil
\end{equation}

\FloatBarrier
\section{Measures of Central Tendency and Dispersion}

\subsection{Central Tendency}
\begin{itemize}
    \item \textbf{Sample Mean:}
    \begin{equation}
        \bar{x} = \frac{1}{m} \sum_{i=1}^m x_i
    \end{equation}
    Highly vulnerable to severe outliers.
    \item \textbf{Median:} Positional center in sorted arrays:
    \begin{equation}
        \text{median} = \begin{cases}
        x_{\left(\frac{m+1}{2}\right)} & \text{if } m \text{ is odd} \\[6pt]
        \dfrac{x_{\left(\frac{m}{2}\right)} + x_{\left(\frac{m}{2} + 1\right)}}{2} & \text{if } m \text{ is even}
        \end{cases}
    \end{equation}
    A robust indicator resistant to extreme values.
    \item \textbf{Mode:} Most frequently observed value; uniquely defined across nominal attributes.
    \item \textbf{Trimmed Mean:} Evaluated by dropping a proportion $\alpha$ of extremes.
\end{itemize}

\subsection{Measures of Dispersion}
\begin{itemize}
    \item \textbf{Range:} $\text{Range} = x_{\max} - x_{\min}$.
    \item \textbf{Sample Variance ($s^2$):}
    \begin{equation}
        s^2 = \frac{1}{m-1} \sum_{i=1}^m (x_i - \bar{x})^2
    \end{equation}
    using Bessel's correction ($m-1$) for unbiased population estimation.
    \item \textbf{Sample Standard Deviation ($s$):} $s = \sqrt{s^2}$.
\end{itemize}

\paragraph{Robust Dispersion Indicators:}
\begin{itemize}
    \item \textbf{Absolute Average Deviation (AAD):}
    \begin{equation}
        \text{AAD}(x) = \frac{1}{m} \sum_{i=1}^m |x_i - \bar{x}|
    \end{equation}
    \item \textbf{Median Absolute Deviation (MAD):}
    \begin{equation}
        \text{MAD}(x) = \text{median}\left(\Big\{|x_i - \text{median}(x)|\Big\}_{i=1}^m\right)
    \end{equation}
\end{itemize}

\FloatBarrier
\section{Five-Number Summary, Quartiles, and Box Plots}

\subsection{The Five-Number Summary}
The **Five-Number Summary** characterizes distributions using positional order statistics:
\begin{equation}
    \Big\{x_{\min},\, Q_1,\, \text{Median } (Q_2),\, Q_3,\, x_{\max}\Big\}
\end{equation}
\begin{itemize}
    \item $Q_1$ (First Quartile): 25th percentile.
    \item $Q_2$ (Median): 50th percentile.
    \item $Q_3$ (Third Quartile): 75th percentile.
    \item \textbf{Interquartile Range (IQR):} $\text{IQR} = Q_3 - Q_1$.
\end{itemize}

\subsection{Box Plot Anatomy}
The **Box Plot** maps the five-number summary geometrically:

\begin{figure}[H]
\centering
\begin{tikzpicture}[scale=0.95]
    % Horizontal axis
    \draw[-Stealth, thick, draw=gray!70] (-5.5, 0) -- (6.5, 0) node[right, font=\small] {$x$};

    % Box from Q1 (-1.2) to Q3 (2.2)
    \fill[navy!12, draw=navy!85!black, thick] (-1.2, 0.7) rectangle (2.2, 2.3);
    
    % Median line at 0.3
    \draw[line width=2pt, draw=navy!95!black] (0.3, 0.7) -- (0.3, 2.3);
    \draw[dotted, navy!60] (0.3, 0.7) -- (0.3, 0.15);

    % IQR brace above
    \draw[decorate, decoration={brace, amplitude=6pt}, thick, draw=navy!85!black] (-1.2, 2.5) -- (2.2, 2.5)
        node[midway, above=6pt, font=\small\bfseries, text=navy] {$\text{IQR} = Q_3 - Q_1$};

    % Left whisker
    \draw[thick, draw=navy!85!black] (-1.2, 1.5) -- (-3.5, 1.5);
    \draw[thick, draw=navy!85!black] (-3.5, 1.1) -- (-3.5, 1.9);

    % Right whisker
    \draw[thick, draw=navy!85!black] (2.2, 1.5) -- (4.8, 1.5);
    \draw[thick, draw=navy!85!black] (4.8, 1.1) -- (4.8, 1.9);

    % Ticks
    \draw[gray!80!black, thick] (-3.5, -0.15) -- (-3.5, 0.15) node[below=3pt, font=\scriptsize\bfseries, align=center, text=crimson] {$L$\\[1pt]\tiny $Q_1 - 1.5\,\text{IQR}$};
    \draw[gray!80!black, thick] (-1.2, -0.15) -- (-1.2, 0.15) node[below=3pt, font=\scriptsize\bfseries] {$Q_1$};
    \draw[navy!80!black, thick] (0.3, -0.15) -- (0.3, 0.15) node[below=3pt, font=\scriptsize\bfseries, align=center, text=navy] {Median\\($Q_2$)};
    \draw[gray!80!black, thick] (2.2, -0.15) -- (2.2, 0.15) node[below=3pt, font=\scriptsize\bfseries] {$Q_3$};
    \draw[gray!80!black, thick] (4.8, -0.15) -- (4.8, 0.15) node[below=3pt, font=\scriptsize\bfseries, align=center, text=crimson] {$U$\\[1pt]\tiny $Q_3 + 1.5\,\text{IQR}$};

    % Outliers
    \fill[crimson] (-4.5, 1.5) circle (3pt) node[above=4pt, font=\scriptsize\bfseries, text=crimson] {Outlier};
    \fill[crimson] (5.6, 1.5) circle (3pt) node[above=4pt, font=\scriptsize\bfseries, text=crimson] {Outlier};
    \fill[crimson] (6.1, 1.5) circle (3pt);
\end{tikzpicture}
\caption{Box Plot anatomy with Tukey fences for univariate outlier demarcation.}
\label{fig:boxplot-anatomy-en}
\end{figure}

\begin{figure}[H]
\centering
\includegraphics[width=0.88\textwidth]{iris_conditional_boxplot.pdf}
\caption{Empirical conditional box plots on the Iris dataset: petal length distributions partitioned by species.}
\label{fig:iris-conditional-boxplot-en}
\end{figure}

\paragraph{Limitations of Box Plots Relative to Histograms:}
Box plots cannot detect multimodality. Identical five-number summaries can emerge from both bimodal and unimodal distributions; histograms remain necessary to reveal intermediate concavities.
